\chapter{Unit 2: Natural Resources, Ecosystem Energetics \& Succession}

\begin{formulabox}{Resource Classifications \& Depletion Parameters}
\begin{itemize}[leftmargin=1.2em]
    \item \textbf{Renewability Criterion}: A resource is renewable if replenishment rate $R_{\text{regen}} \ge \text{Harvest rate } H_{\text{draw}}$. If $H_{\text{draw}} > R_{\text{regen}}$, resource exhaustion occurs.
    \item \textbf{Soil Formation Rate}: Pedogenesis generates approximately $1\text{ cm}$ of fertile topsoil every $200\text{--}400\text{ years}$.
    \item \textbf{Universal Soil Loss Equation (USLE)}:
    \begin{equation}
    A = R \cdot K \cdot LS \cdot C \cdot P
    \end{equation}
    where $A = \text{computed soil loss}$, $R = \text{rainfall erosivity factor}$, $K = \text{soil erodibility factor}$, $LS = \text{slope length-gradient factor}$, $C = \text{crop cover factor}$, $P = \text{conservation practice factor}$.
    \item \textbf{Coastal Aquifer Saltwater Intrusion (Ghyben-Herzberg Relation)}:
    \begin{equation}
    z = \frac{\rho_f}{\rho_s - \rho_f} h \approx 40 h
    \end{equation}
    where $h = \text{freshwater head above sea level}$, $z = \text{depth to saltwater interface below sea level}$, $\rho_f = 1.000\text{ g/cm}^3$, $\rho_s = 1.025\text{ g/cm}^3$. Every $1\text{ meter}$ drop in groundwater table induces a $40\text{ meter}$ saltwater rise.
    \item \textbf{National Forest Policy Targets (1988)}:
    \begin{itemize}
        \item National target: Minimum $33\%$ of total geographical land area under forest and tree cover.
        \item Hill and mountainous ecosystems: Minimum $67\%$ ($2/3\text{rd}$) forest cover for slope stability and soil conservation.
    \end{itemize}
    \item \textbf{Wind Power Available Kinetic Energy}:
    \begin{equation}
    P_{\text{wind}} = \frac{1}{2}\rho A v^3
    \end{equation}
    where $\rho = \text{air density } (\approx 1.225\text{ kg/m}^3)$, $A = \text{swept rotor area } (\pi R^2)$, $v = \text{wind velocity}$. Power scales with the cube of wind velocity ($v^3$). Betz limit maximum aerodynamic efficiency $= 59.3\%$ ($16/27$).
\end{itemize}
\end{formulabox}

\begin{formulabox}{Ecosystem Energetics \& Trophic Transfer Efficiencies}
\begin{itemize}[leftmargin=1.2em]
    \item \textbf{Primary Productivity Formulations}:
    \begin{align}
    \text{GPP} &= \text{Total rate of organic chemical energy synthesized by autotrophs via photosynthesis} \\
    \text{NPP} &= \text{GPP} - R_a
    \end{align}
    where $R_a = \text{autotrophic respiration energy expended for cellular maintenance}$.
    \item \textbf{Net Ecosystem Productivity (NEP)}:
    \begin{equation}
    \text{NEP} = \text{GPP} - (R_a + R_h) = \text{NPP} - R_h
    \end{equation}
    where $R_h = \text{heterotrophic respiration (consumers + decomposers)}$.
    \item \textbf{Lindeman's 10\% Trophic Efficiency Rule (1942)}:
    \begin{equation}
    \text{Ecological Efficiency } (\lambda) = \frac{E_n}{E_{n-1}} \times 100\% \approx 10\%
    \end{equation}
    Approximately $90\%$ of energy is lost at each trophic link via metabolic respiration ($R$), unconsumed biomass, and heat dissipation according to the Second Law of Thermodynamics.
    \item \textbf{Ecological Pyramids Comparison}:
    \begin{itemize}
        \item \textbf{Pyramid of Numbers}: Upright in grasslands; inverted in parasitic food chains; spindle-shaped in temperate forest trees.
        \item \textbf{Pyramid of Biomass}: Upright in terrestrial forests; \textbf{inverted in marine open waters} (phytoplankton standing crop is smaller than zooplankton due to rapid turnover).
        \item \textbf{Pyramid of Energy}: \textbf{ALWAYS UPRIGHT} in all ecosystems without exception.
    \end{itemize}
\end{itemize}
\end{formulabox}

\begin{formulabox}{Ecological Succession Dynamics}
\begin{itemize}[leftmargin=1.2em]
    \item \textbf{Primary vs Secondary Succession}:
    \begin{itemize}
        \item \textbf{Primary Succession}: Initiates on sterile, virgin substrates lacking soil (bare rock, cooled volcanic basalt, sand dunes, retreating glaciers). Pioneers: Lichens, mosses.
        \item \textbf{Secondary Succession}: Occurs on disturbed sites where pre-existing soil and seed banks remain (abandoned farmland, burned forest). Significantly faster rate.
    \end{itemize}
    \item \textbf{Six Seral Stages of Succession (Clementsian Model)}:
    \begin{enumerate}
        \item \textbf{Nudation}: Development of bare exposed substrate.
        \item \textbf{Invasion / Migration}: Dispersal and arrival of propagules (seeds, spores).
        \item \textbf{Ecesis}: Successful establishment, germination, and growth of pioneers.
        \item \textbf{Aggregation}: Population increase and clustering of established species.
        \item \textbf{Competition}: Interspecific and intraspecific struggle for space, nutrients, and light.
        \item \textbf{Reaction \& Climax Stabilization}: Biotic modification of soil and microclimate, culminating in a stable self-perpetuating climax community.
    \end{enumerate}
    \item \textbf{Trends from Pioneer to Climax}:
    \begin{itemize}
        \item Gross production / community respiration ($P/R$ ratio): $P/R > 1$ in early seral stages; $P/R \to 1.0$ at climax equilibrium.
        \item Biomass: Low $\to$ High. \quad Diversity ($H'$): Low $\to$ High. \quad Food webs: Linear chains $\to$ Complex webs.
    \end{itemize}
\end{itemize}
\end{formulabox}

\begin{memorytip}
The pyramid of energy is NEVER inverted! In open oceans, the biomass pyramid is inverted because phytoplankton reproduce and die in days, but the energy throughput per year is always higher at the base!
\end{memorytip}
